\documentclass[10pt,a4paper]{amsart}
\usepackage{amsmath,amssymb,amsthm,graphicx}
\usepackage{verbatim}
\newenvironment{reference}{\comment}{\endcomment}
%\newenvironment{reference}{}{}
\newenvironment{slogan}{\comment}{\endcomment}
\newenvironment{history}{\comment}{\endcomment}

% For commutative diagrams we use Xy-pic
\usepackage[all]{xy}

% We use 2cell for 2-commutative diagrams.
\xyoption{2cell}
% Theorem environments.
%
\theoremstyle{plain}
\newtheorem{theorem}[subsection]{Theorem}
\newtheorem{proposition}[subsection]{Proposition}
\newtheorem{lemma}[subsection]{Lemma}

\theoremstyle{definition}
\newtheorem{definition}[subsection]{Definition}
\newtheorem{example}[subsection]{Example}
\newtheorem{exercise}[subsection]{Exercise}
\newtheorem{situation}[subsection]{Situation}

\theoremstyle{remark}
\newtheorem{remark}[subsection]{Remark}
\newtheorem{remarks}[subsection]{Remarks}

\numberwithin{equation}{subsection}

% Macros
%
\def\lim{\mathop{\mathrm{lim}}\nolimits}
\def\colim{\mathop{\mathrm{colim}}\nolimits}
\def\Spec{\mathop{\mathrm{Spec}}}
\def\Hom{\mathop{\mathrm{Hom}}\nolimits}
\def\Ext{\mathop{\mathrm{Ext}}\nolimits}
\def\SheafHom{\mathop{\mathcal{H}\!\mathit{om}}\nolimits}
\def\SheafExt{\mathop{\mathcal{E}\!\mathit{xt}}\nolimits}
\def\Sch{\mathit{Sch}}
\def\Mor{\mathop{\mathrm{Mor}}\nolimits}
\def\Ob{\mathop{\mathrm{Ob}}\nolimits}
\def\Sh{\mathop{\mathit{Sh}}\nolimits}
\def\NL{\mathop{N\!L}\nolimits}
\def\CH{\mathop{\mathrm{CH}}\nolimits}
\def\proetale{{pro\text{-}\acute{e}tale}}
\def\etale{{\acute{e}tale}}
\def\QCoh{\mathit{QCoh}}
\def\Ker{\mathop{\mathrm{Ker}}}
\def\Im{\mathop{\mathrm{Im}}}
\def\Coker{\mathop{\mathrm{Coker}}}
\def\Coim{\mathop{\mathrm{Coim}}}

% Boxtimes
%
\DeclareMathSymbol{\boxtimes}{\mathbin}{AMSa}{"02}

%
% Macros for moduli stacks/spaces
%
\def\QCohstack{\mathcal{QC}\!\mathit{oh}}
\def\Cohstack{\mathcal{C}\!\mathit{oh}}
\def\Spacesstack{\mathcal{S}\!\mathit{paces}}
\def\Quotfunctor{\mathrm{Quot}}
\def\Hilbfunctor{\mathrm{Hilb}}
\def\Curvesstack{\mathcal{C}\!\mathit{urves}}
\def\Polarizedstack{\mathcal{P}\!\mathit{olarized}}
\def\Complexesstack{\mathcal{C}\!\mathit{omplexes}}
% \Pic is the operator that assigns to X its picard group, usage \Pic(X)
% \Picardstack_{X/B} denotes the Picard stack of X over B
% \Picardfunctor_{X/B} denotes the Picard functor of X over B
\def\Pic{\mathop{\mathrm{Pic}}\nolimits}
\def\Picardstack{\mathcal{P}\!\mathit{ic}}
\def\Picardfunctor{\mathrm{Pic}}
\def\Deformationcategory{\mathcal{D}\!\mathit{ef}}
\usepackage{cleveref}
\crefname{originalreference}{original reference}{original references}
\Crefname{originalreference}{Original reference}{Original references}
\makeatletter
\newlabel{varieties-lemma-algebraic-scheme-dim-0}{{06LH}{source}{}{sourcepointer-1}{}}
\newlabel{varieties-lemma-algebraic-scheme-dim-0@cref}{{[originalreference][0][]06LH}{[1][0][]source}{}{}{}}
\makeatother
\crefname{originalreference}{original reference}{original references}
\Crefname{originalreference}{Original reference}{Original references}
\makeatletter
\newlabel{example-general-linear-group}{{022W}{source}{}{sourcepointer-2}{}}
\newlabel{example-general-linear-group@cref}{{[originalreference][0][]022W}{[1][0][]source}{}{}{}}
\makeatother
\crefname{originalreference}{original reference}{original references}
\Crefname{originalreference}{Original reference}{Original references}
\makeatletter
\newlabel{algebra-remark-intersection-powers-ideal}{{00IR}{source}{}{sourcepointer-3}{}}
\newlabel{algebra-remark-intersection-powers-ideal@cref}{{[originalreference][0][]00IR}{[1][0][]source}{}{}{}}
\makeatother
\crefname{originalproposition}{original proposition}{original propositions}
\Crefname{originalproposition}{Original proposition}{Original propositions}
\makeatletter
\newlabel{proposition-connected-component}{{0B7R}{source}{}{sourcepointer-4}{}}
\newlabel{proposition-connected-component@cref}{{[originalproposition][0][]0B7R}{[1][0][]source}{}{}{}}
\makeatother
\crefname{originallemma}{original lemma}{original lemmas}
\Crefname{originallemma}{Original lemma}{Original lemmas}
\makeatletter
\newlabel{lemma-irreducible-group-scheme-over-field-qc}{{0B7P}{source}{}{sourcepointer-5}{}}
\newlabel{lemma-irreducible-group-scheme-over-field-qc@cref}{{[originallemma][0][]0B7P}{[1][0][]source}{}{}{}}
\makeatother
\crefname{originalreference}{original reference}{original references}
\Crefname{originalreference}{Original reference}{Original references}
\makeatletter
\newlabel{descent-lemma-closed-immersion}{{03I0}{source}{}{sourcepointer-6}{}}
\newlabel{descent-lemma-closed-immersion@cref}{{[originalreference][0][]03I0}{[1][0][]source}{}{}{}}
\makeatother
\crefname{originalreference}{original reference}{original references}
\Crefname{originalreference}{Original reference}{Original references}
\makeatletter
\newlabel{descent-lemma-descending-property-closed-immersion}{{02L6}{source}{}{sourcepointer-7}{}}
\newlabel{descent-lemma-descending-property-closed-immersion@cref}{{[originalreference][0][]02L6}{[1][0][]source}{}{}{}}
\makeatother
\title{The center of a locally algebraic group scheme is a closed subgroup scheme}
\author{The Stacks Project Authors}
\date{stacks-project:groupoids.tex:local-snapshot-20261012; GNU FDL 1.2 or later}
\begin{document}
\maketitle
\noindent\textit{Editorial note.} All delivered mathematical text is the author's own, copied verbatim from the pinned Stacks Project chapter file groupoids.tex; no statement or proof step was rewritten, reordered or removed.\par\medskip
\noindent\textit{Editorial note.} Span recomputation by label pairing (LANE\_BRIEF 10.2): in this chapter the author's proof environment follows the statement environment as a standalone block; the spec covers the complete environment pair (statement span, proof span) as recorded in the pool candidate.\par\medskip
\noindent\textit{Editorial note.} Occupation scan: this chapter's pool has 47 of 128 labels already delivered by the S7 line (S7-gp-* registrations, statement\_anchor-prose scan; withdrawn cards excluded); none of the 11 accepted labels of this batch is among them (gp-occupied.json, 2026-10-14).\par\medskip
\ifdefined\hypertarget\hypertarget{sourcepointer-1}{}\fi
\noindent\textit{Source reference.} Original reference 06LH: sources/stacks-chow/tags.txt, source lines 8466--8466. Cross-chapter standard prerequisite cited by the author inside the counted span; it is used as a pointer and its own proof is not part of this item.\par\medskip
\ifdefined\hypertarget\hypertarget{sourcepointer-2}{}\fi
\noindent\textit{Source reference.} Original reference 022W: sources/stacks-chow/tags.txt, source lines 2721--2721. Cross-chapter standard prerequisite cited by the author inside the counted span; it is used as a pointer and its own proof is not part of this item.\par\medskip
\ifdefined\hypertarget\hypertarget{sourcepointer-3}{}\fi
\noindent\textit{Source reference.} Original reference 00IR: sources/stacks-chow/tags.txt, source lines 689--689. Cross-chapter standard prerequisite cited by the author inside the counted span; it is used as a pointer and its own proof is not part of this item.\par\medskip
\ifdefined\hypertarget\hypertarget{sourcepointer-4}{}\fi
\noindent\textit{Source reference.} Original proposition 0B7R: sources/stacks-chow/tags.txt, source lines 14114--14114. Cross-chapter standard prerequisite cited by the author inside the counted span; it is used as a pointer and its own proof is not part of this item.\par\medskip
\ifdefined\hypertarget\hypertarget{sourcepointer-5}{}\fi
\noindent\textit{Source reference.} Original lemma 0B7P: sources/stacks-chow/tags.txt, source lines 14112--14112. Cross-chapter standard prerequisite cited by the author inside the counted span; it is used as a pointer and its own proof is not part of this item.\par\medskip
\ifdefined\hypertarget\hypertarget{sourcepointer-6}{}\fi
\noindent\textit{Source reference.} Original reference 03I0: sources/stacks-chow/tags.txt, source lines 4586--4586. Cross-chapter standard prerequisite cited by the author inside the counted span; it is used as a pointer and its own proof is not part of this item.\par\medskip
\ifdefined\hypertarget\hypertarget{sourcepointer-7}{}\fi
\noindent\textit{Source reference.} Original reference 02L6: sources/stacks-chow/tags.txt, source lines 3397--3397. Cross-chapter standard prerequisite cited by the author inside the counted span; it is used as a pointer and its own proof is not part of this item.\par\medskip
\begin{lemma}
\label{lemma-algebraic-center}
Let $k$ be a field. Let $G$ be a locally algebraic group scheme over $k$.
Then the center of $G$ is a closed subgroup scheme of $G$.
\end{lemma}

\begin{proof}
Let $\text{Aut}(G)$ denote the contravariant functor on the category of
schemes over $k$ which associates to $S/k$ the set of automorphisms
of the base change $G_S$ as a group scheme over $S$. There is a natural
transformation
$$
G \longrightarrow \text{Aut}(G),\quad
g \longmapsto \text{inn}_g
$$
sending an $S$-valued point $g$ of $G$ to the inner automorphism of $G$
determined by $g$. The center $C$ of $G$ is by definition the kernel of
this transformation, i.e., the functor which to $S$ associates those
$g \in G(S)$ whose associated inner automorphism is trivial. The statement
of the lemma is that this functor is representable by a closed subgroup
scheme of $G$.

\medskip\noindent
Choose an integer $n \geq 1$. Let $G_n \subset G$ be the $n$th infinitesimal
neighbourhood of the identity element $e$ of $G$. For every scheme $S/k$
the base change $G_{n, S}$ is the $n$th infinitesimal neighbourhood of
$e_S : S \to G_S$. Thus we see that there is a natural transformation
$\text{Aut}(G) \to \text{Aut}(G_n)$ where the right hand side is the
functor of automorphisms of $G_n$ as a scheme ($G_n$ isn't in general
a group scheme). Observe that $G_n$ is the spectrum of an artinian
local ring $A_n$ with residue field $k$ which has finite dimension
as a $k$-vector space
(Varieties, Lemma \ref{varieties-lemma-algebraic-scheme-dim-0}).
Since every automorphism of $G_n$ induces in particular an invertible
linear map $A_n \to A_n$, we obtain transformations of functors
$$
G \to \text{Aut}(G) \to \text{Aut}(G_n) \to \text{GL}(A_n)
$$
The final group valued functor is representable, see
Example \ref{example-general-linear-group}, and the
last arrow is visibly injective.
Thus for every $n$ we obtain a closed subgroup scheme
$$
H_n = \Ker(G \to \text{Aut}(G_n)) = \Ker(G \to \text{GL}(A_n)).
$$
As a first approximation we set $H = \bigcap_{n \geq 1} H_n$
(scheme theoretic intersection). This is a closed subgroup scheme
which contains the center $C$.

\medskip\noindent
Let $h$ be an $S$-valued point of $H$ with $S$ locally Noetherian.
Then the automorphism $\text{inn}_h$ induces the identity on all
the closed subschemes $G_{n, S}$. Consider the kernel
$K = \Ker(\text{inn}_h : G_S \to G_S)$.
This is a closed subgroup scheme of $G_S$ over $S$ containing the
closed subschemes $G_{n, S}$ for $n \geq 1$.
This implies that $K$ contains an open neighbourhood of
$e(S) \subset G_S$, see
Algebra, Remark \ref{algebra-remark-intersection-powers-ideal}.
Let $G^0 \subset G$ be as in Proposition \ref{proposition-connected-component}.
Since $G^0$ is geometrically irreducible, we conclude that $K$ contains
$G^0_S$ (for any nonempty open $U \subset G^0_{k'}$ and any field extension
$k'/k$ we have $U \cdot U^{-1} = G^0_{k'}$, see proof of
Lemma \ref{lemma-irreducible-group-scheme-over-field-qc}).
Applying this with $S = H$ we find that $G^0$ and $H$
are subgroup schemes of $G$ whose points commute: for any scheme $S$
and any $S$-valued points $g \in G^0(S)$, $h \in H(S)$ we have
$gh = hg$ in $G(S)$.

\medskip\noindent
Assume that $k$ is algebraically closed. Then we can pick a $k$-valued
point $g_i$ in each irreducible component $G_i$ of $G$. Observe that in
this case the connected components of $G$ are the irreducible components
of $G$ are the translates of $G^0$ by our $g_i$. We claim that
$$
C = H \cap \bigcap\nolimits_i \Ker(\text{inn}_{g_i} : G \to G)
\quad (\text{scheme theoretic intersection})
$$
Namely, $C$ is contained in the right hand side. On the other hand, every
$S$-valued point $h$ of the right hand side commutes with $G^0$
and with $g_i$ hence with everything in $G = \bigcup G^0g_i$.

\medskip\noindent
The case of a general base field $k$ follows from the result for the
algebraic closure $\overline{k}$ by descent. Namely, let
$A \subset G_{\overline{k}}$ the closed subgroup scheme representing
the center of $G_{\overline{k}}$. Then we have
$$
A \times_{\Spec(k)} \Spec(\overline{k}) =
\Spec(\overline{k}) \times_{\Spec(k)} A
$$
as closed subschemes of $G_{\overline{k} \otimes_k \overline{k}}$ by
the functorial nature of the center. Hence we see that $A$ descends
to a closed subgroup scheme $Z \subset G$ by
Descent, Lemma \ref{descent-lemma-closed-immersion}
(and Descent, Lemma \ref{descent-lemma-descending-property-closed-immersion}).
Then $Z$ represents $C$ (small argument omitted) and the proof is complete.
\end{proof}

\end{document}
